% from ega2-1-fr.tex
\oldpage[II]{57}

et en particulier
\[
\label{II.3.3.1.2-fr}
  \Gamma_*(\mathcal{O}_X)
  =\bigoplus_{n\in\mathbf{Z}}p_*(\mathcal{O}_X(n)).
\tag{3.3.1.2}
\]
On sait (\hyperref[0.4.2.2-fr]{0, 4.2.2}) qu'il existe un homomorphisme
canonique
\[
  p_*(\mathcal{F})\otimes_{\mathcal{O}_Y}p_*(\mathcal{G})
  \longrightarrow
  p_*(\mathcal{F}\otimes_{\mathcal{O}_X}\mathcal{G})
\]
pour deux $\mathcal{O}_X$-Modules $\mathcal{F}$ et $\mathcal{G}$~; on déduit
donc de (\hyperref[II.3.2.7.1-fr]{3.2.7.1}) que
$\Gamma_*(\mathcal{O}_X)$ est muni d'une structure de
$\mathcal{O}_Y$-\emph{Algèbre graduée}, et
(\hyperref[II.3.2.5.2-fr]{3.2.5.2}) définit de même sur
$\Gamma_*(\mathcal{F})$ une structure de \emph{Module gradué} sur
$\Gamma_*(\mathcal{O}_X)$.

En vertu de (\hyperref[II.3.2.5-fr]{3.2.5}), et de l'exactitude à gauche du
foncteur $p_*$ (\hyperref[0.4.2.1-fr]{0, 4.2.1}),
$\Gamma_*(\mathcal{F})$ est un foncteur covariant additif et
\emph{exact à gauche} en $\mathcal{F}$ dans la catégorie des
$\mathcal{O}_X$-Modules, à valeurs dans la catégorie des
$\mathcal{O}_Y$-Modules gradués (où les morphismes sont les homomorphismes de
degré $0$). En particulier, si $\mathcal{J}$ est un faisceau d'idéaux dans
$\mathcal{O}_X$, $\Gamma_*(\mathcal{J})$ s'identifie à un
\emph{faisceau gradué d'idéaux} de $\Gamma_*(\mathcal{O}_X)$.
\end{env}

\begin{env}[3.3.2]
\label{II.3.3.2-fr}
Soit $\mathcal{M}$ un $\mathcal{S}$-Module gradué quasi-cohérent. Pour tout
ouvert affine $U$ de $Y$, on a défini en
(\hyperref[II.2.6.2-fr]{2.6.2}) un homomorphisme de groupes abéliens
\[
  \alpha_{0,U}:\Gamma(U,\mathcal{M}_0)
  \longrightarrow\Gamma(p^{-1}(U),\widetilde{\mathcal{M}}).
\]
Il est immédiat que ces homomorphismes commutent aux opérations de restriction
(\hyperref[II.2.8.13.1-fr]{2.8.13.1}) et définissent donc (sans utiliser
l'hypothèse que $\mathcal{S}$ est engendrée par $\mathcal{S}_1$) un
homomorphisme de faisceaux de groupes abéliens
\[
\label{II.3.3.2.1-fr}
  \alpha_0:\mathcal{M}_0\longrightarrow p_*(\widetilde{\mathcal{M}}).
\tag{3.3.2.1}
\]
Appliquant ce résultat à chacun des
$\mathcal{M}_n=(\mathcal{M}(n))_0$, et tenant compte de
(\hyperref[II.3.2.8.1-fr]{3.2.8.1}), on définit un homomorphisme de faisceaux
de groupes abéliens
\[
\label{II.3.3.2.2-fr}
  \alpha_n:\mathcal{M}_n\longrightarrow
  p_*(\widetilde{\mathcal{M}}(n))
  \qquad\text{pour tout }n\in\mathbf{Z},
\tag{3.3.2.2}
\]
d'où un homomorphisme fonctoriel (de degré $0$) de faisceaux gradués de groupes
abéliens
\[
\label{II.3.3.2.3-fr}
  \alpha:\mathcal{M}\longrightarrow\Gamma_*(\widetilde{\mathcal{M}})
\tag{3.3.2.3}
\]
(aussi noté $\alpha_{\mathcal{M}}$).

En prenant en particulier $\mathcal{M}=\mathcal{S}$, on vérifie que
$\alpha:\mathcal{S}\to\Gamma_*(\mathcal{O}_X)$ est un homomorphisme de
$\mathcal{O}_Y$-Algèbres graduées et que
(\hyperref[II.3.3.2.3-fr]{3.3.2.3}) est un di-homomorphisme de Modules
gradués, relatif à cet homomorphisme d'Algèbres graduées.

Remarquons encore qu'à chacun des $\alpha_n$ il correspond
(\hyperref[0.4.4.3-fr]{0, 4.4.3}) un homomorphisme canonique de
$\mathcal{O}_X$-Modules
\[
\label{II.3.3.2.4-fr}
  \alpha_n^\sharp:p^*(\mathcal{M}_n)
  \longrightarrow\widetilde{\mathcal{M}}(n).
\tag{3.3.2.4}
\]
On vérifie sans peine que cet homomorphisme n'est autre que celui qui
correspond fonctoriellement (\hyperref[II.3.2.4-fr]{3.2.4}) à
l'homomorphisme canonique (de degré $0$) de $\mathcal{O}_Y$-Modules gradués
\[
\label{II.3.3.2.5-fr}
  \mathcal{M}_n\otimes_{\mathcal{O}_Y}\mathcal{S}
  \longrightarrow\mathcal{M}(n)
\tag{3.3.2.5}
\]

